\documentclass{article}
\usepackage{amsmath,amssymb}
\usepackage{xurl}
\usepackage[hidelinks]{hyperref}
% Shared commands for notes in docs/paper-gaps/.

\DeclareUnicodeCharacter{2080}{\ifmmode{}_{0}\else\textsubscript{0}\fi}
\DeclareUnicodeCharacter{2081}{\ifmmode{}_{1}\else\textsubscript{1}\fi}
\DeclareUnicodeCharacter{2082}{\ifmmode{}_{2}\else\textsubscript{2}\fi}
\DeclareUnicodeCharacter{2099}{\ifmmode{}_{n}\else\textsubscript{n}\fi}

\newcommand{\Lean}{\textsc{Lean}}
\ExplSyntaxOn
\NewDocumentCommand{\leanid}{m}{
  \ifmmode
    \textsf{\detokenize{#1}}
  \else
    \group_begin:
    \urlstyle{sf}
    \tl_set:Nn \l_tmpa_tl {#1}
    \tl_replace_all:Nnn \l_tmpa_tl {\_} {_}
    \exp_args:NV \path \l_tmpa_tl
    \group_end:
  \fi
}
\ExplSyntaxOff
\providecommand{\tr}{\operatorname{tr}}
\newcommand{\tnleanIssueBase}{https://github.com/LionSR/TNLean/issues/}
\newcommand{\tnleanPrBase}{https://github.com/LionSR/TNLean/pull/}
\newcommand{\tnleanDocBase}{https://sirui-lu.com/TNLean/docs/}
\newcommand{\ghissue}[1]{\href{\tnleanIssueBase#1}{GitHub issue \##1}}
\newcommand{\ghpr}[1]{\href{\tnleanPrBase#1}{PR \##1}}
\newcommand{\doclink}[1]{\href{\tnleanDocBase#1.html}{\path{#1}}}

% Dirac notation and the identity operator, as in blueprint/src/macros/common.tex.
% \providecommand keeps note-local definitions authoritative.
\usepackage{dsfont}
\providecommand{\ket}[1]{|#1\rangle}
\providecommand{\bra}[1]{\langle #1|}
\providecommand{\braket}[2]{\langle #1 | #2 \rangle}
\providecommand{\ketbra}[2]{|#1\rangle\!\langle #2|}
\providecommand{\Id}{\mathds{1}}
\providecommand{\one}{\mathds{1}}
\providecommand{\C}{\mathbb{C}}
\providecommand{\MN}[1]{M_{#1}(\C)}
\providecommand{\MD}{M_D(\C)}

% External referencing for paper sources.  The build helper writes sanitized
% auxiliary files under build/paper-aux/ and excludes duplicated source labels.
\usepackage{xr}

% Import a generated paper auxiliary file with a caller-supplied label prefix.
% The candidate paths support builds from docs/paper-gaps/, the repository root,
% and build/paper-aux/ itself.
\newcommand{\importpaperaux}[2]{%
  \IfFileExists{../../build/paper-aux/#2.aux}{%
    \externaldocument[#1]{../../build/paper-aux/#2}%
  }{%
    \IfFileExists{build/paper-aux/#2.aux}{%
      \externaldocument[#1]{build/paper-aux/#2}%
    }{%
      \IfFileExists{#2.aux}{%
        \externaldocument[#1]{#2}%
      }{}%
    }%
  }%
}

% General external reference macros
\newcommand{\paperref}[2]{\ref{#1-#2}}
\newcommand{\papereqref}[2]{(\ref{#1-#2})}

% CPSV16 source (1606.00608 / MPDO-22-12-17-2.tex) external document and shortcuts
\importpaperaux{cpsv16-}{1606.00608/MPDO-22-12-17-2-tnlean}
\newcommand{\cpsvref}[1]{\paperref{cpsv16}{#1}}
\newcommand{\cpsveqref}[1]{\papereqref{cpsv16}{#1}}

% Verdict marker: \gapnote{<kind>}{<status>}. Kind is the policy
% classification (clarification, local-correction, scope-restriction,
% unfaithful, false-source, open-gap); status is open, wip (elimination
% underway), resolved, or historical. Machine-read by the site index;
% prints a verdict line.
\newcommand{\gapnote}[2]{\par\noindent\textbf{Verdict:} #1 (#2).\par}

\title{The Positive Part of a Blocked Tensor with Overlapping Blocks}
\author{}
\date{2026-09-26}
\begin{document}
\maketitle
\gapnote{false-source}{open}

\section{What the source states}
Let $A^i=\bigoplus_{j=1}^b\operatorname{diag}(\mu_{j,1},\dots,\mu_{j,m_j})
    \otimes A_j^i$ with normal $A_j$ forming a basis of normal tensors
\cite[eq.~(S2)]{Malz2024Preparation}, block $q$ sites, and write the polar
decomposition of the blocked map as $B=VP$. The source argues ``Since $V$ is
an isometry, $P$ inherits the block structure of $A$'',
\begin{align}
    P & =\bigoplus_{j=1}^b\operatorname{diag}(\mu_{j,1}^q,\dots,\mu_{j,m_j}^q)
    \otimes P_j
    \label{eq:mswc24_mixed_block_form}
\end{align}
with normal $P_j$ \cite[eq.~(S5)]{Malz2024Preparation}. It then approximates the
state on $N=qM$ sites by
$\ket{\widetilde\phi_N}\propto V^{\otimes M}\sum_j\beta_j\ket{\Omega_j}$ with
$\beta_j=\sum_k\mu_{j,k}^N$ \cite[eq.~(S7)]{Malz2024Preparation}, and part~(ii)
of the approximation-error lemma bounds
$\varepsilon=1-\lvert\braket{\widetilde\phi_N}{\phi_N}\rvert$ by
$O\bigl((N/q)e^{-\gamma q/\xi_{\mathrm{diag}}}\bigr)$ for $q=o(N)$
\cite[eq.~(S12)]{Malz2024Preparation}, where
$\xi_{\mathrm{diag}}=\max_j\xi_{jj}$ is the largest correlation length of the
blocks. In the proof, the overlaps of distinct blocks enter only through a
term $O(e^{-N/\xi_{\mathrm{off\text{-}diag}}})$, which $q=o(N)$ removes.

\section{Where the argument fails}
The positive part is $P=(B^\dagger B)^{1/2}$, and $B^\dagger B$ is the transfer
matrix of $B$ with its legs regrouped. Its block between two distinct blocks
$A_j$ and $A_{j'}$ is the $q$-th power of the mixed transfer matrix
$E_{jj'}=\sum_iA_j^{i*}\otimes A_{j'}^i$. The source uses that
$E_{jj'}$ has spectral radius $\tau<1$ \cite[proof of Lemma~1$'$]{Malz2024Preparation},
but $E_{jj'}^q$ is not zero, so $B^\dagger B$ and its square root are not block
diagonal. The block form (\ref{eq:mswc24_mixed_block_form}) holds only when the
$q$-site states of distinct blocks are orthogonal. This failure is separate from
the one for repeated blocks \cite{gap:mswc24_multiplicity_fixed_point}: it occurs
when every multiplicity is one.

\section{The counterexample}
Take $d=2$ and the two one-dimensional normal blocks $A_1=(1,0)$ and
$A_2=(3/5,4/5)$, each of multiplicity one with weight one, so that
$A^0=\operatorname{diag}(1,3/5)$ and $A^1=\operatorname{diag}(0,4/5)$. Both
blocks satisfy $\sum_i\lvert a_i\rvert^2=1$ and are not proportional, so they form a
basis of normal tensors, and $\beta=(1,1)$. The $q$-site states of the blocks are
$\ket{0\cdots0}$ and $(\tfrac35\ket0+\tfrac45\ket1)^{\otimes q}$, with overlap
$s=(3/5)^q$. On the diagonal bond coordinates $(0,0)$ and $(1,1)$, the only ones
on which $B$ does not vanish,
\begin{align}
    B^\dagger B=\begin{pmatrix}1&s\\s&1\end{pmatrix},\qquad
    P=\begin{pmatrix}u&v\\v&u\end{pmatrix},\qquad
    u=\frac{\sqrt{1+s}+\sqrt{1-s}}2,\quad
    v=\frac{\sqrt{1+s}-\sqrt{1-s}}2,
    \label{eq:mswc24_mixed_positive_part}
\end{align}
and the support projector of $P$ is the projector onto these two coordinates. The
fixed point of a one-dimensional block is $\sigma=1$, so the pairs are
$\omega_1=e_{00}$ and $\omega_2=e_{11}$. Since $V^\dagger B=P$ and
$V^\dagger V$ fixes the pairs, for every $q,M\geq1$,
\begin{align}
    \braket{\widetilde\phi_N}{\phi_N}
     & =\frac{u^M+v^M}{(1+s^M)^{1/2}}.
    \label{eq:mswc24_mixed_overlap}
\end{align}
For $q=1$ this is $(3^M+1)/(10^M+6^M)^{1/2}$, which tends to zero. Since
$u^2+v^2=1$, $u^2=(1+\sqrt{1-s^2})/2\geq1/2$, and $1-u\geq s^2/8$, one has
$u^M+v^M\leq u^3+v^2=1-u^2(1-u)\leq1-s^2/16$ for $M\geq3$, so
\begin{align}
    \varepsilon & \geq\frac{s^2}{16}=\frac1{16}\Bigl(\frac9{25}\Bigr)^q
    \qquad(M\geq3).
    \label{eq:mswc24_mixed_lower_bound}
\end{align}

The transfer maps of one-dimensional blocks are the identity and have no
eigenvalue other than one. The source defines
$\xi_{jj}=-1/\ln\lvert\lambda_2^{(j)}\rvert$ with $\lambda_2^{(j)}$ the subleading
eigenvalue of the transfer matrix of $A_j$ \cite{Malz2024Preparation}; here there
is none, so $\xi_{\mathrm{diag}}$ is degenerate. Under the usual convention
$\lambda_2^{(j)}=0$ it is $0$, and the printed bound then asserts $\varepsilon=0$,
which (\ref{eq:mswc24_mixed_lower_bound}) refutes directly. The formal statement
avoids the convention: it reads $\xi_{\mathrm{diag}}$ as the correlation length at any
$\lambda_2$ with $0<\lvert\lambda_2\rvert<1$ bounding the subleading eigenvalues,
which here is every such $\lambda_2$, and covers those with
$r=e^{-\gamma/\xi_{\mathrm{diag}}}=\lvert\lambda_2\rvert^\gamma<9/25$. Take the
sequence $q=M$, $N=q^2$, which satisfies $q=o(N)$.
The printed bound $C(N/q)e^{-\gamma q/\xi_{\mathrm{diag}}}=Cqr^q$ is eventually
smaller than $(9/25)^q/16$ for every constant $C$, and so is
$Cqr^q\exp(Cqr^q)$. The error is governed by the overlap $s=\tau^q$ of the blocks,
with $\tau=3/5$ the eigenvalue of the mixed transfer matrix, and is of order
$Ms^2=(N/q)\tau^{2q}$, a contribution that the source's estimate does not contain.

\section{What is formalized}
The overlap (\ref{eq:mswc24_mixed_overlap}) for all $q,M\geq1$, the lower bound
(\ref{eq:mswc24_mixed_lower_bound}), the fact that every $\lambda_2$ bounds the
subleading eigenvalues of the two blocks, and the failure of the printed bound
along $q=M$ are formalized.\footnote{Formal declarations:
    \leanid{MPSTensor.nonNormalApproxOverlap_overlappingBlockTensor},
    \leanid{MPSTensor.le_one_sub_norm_nonNormalApproxOverlap_overlappingBlockTensor},
    \leanid{MPSTensor.norm_le_of_hasEigenvalue_transferMap_overlappingBlockBasis},
    \leanid{MPSTensor.not_approximationError_le_overlappingBlockTensor} in
    \doclink{TNLean/MPS/Preparation/OverlappingBlockCounterexample}.} The
positive part and the support projector are computed from the uniqueness of the
positive square root of $B^\dagger B$.

\section{Elimination}
Two statements would repair part~(ii) for multiplicity one. The first adds the
hypothesis that the $q$-site states of distinct blocks are orthogonal,
$\sum_{i_1,\dots,i_q}\overline{(A_j^{i_1}\cdots A_j^{i_q})_{ab}}\,
(A_{j'}^{i_1}\cdots A_{j'}^{i_q})_{a'b'}=0$ for $j\neq j'$. Under it
(\ref{eq:mswc24_mixed_block_form}) holds with $\lvert\mu_j\rvert^q$ in place of
$\mu_j^q$: writing the blocked map as $B=\sum_j\mu_j^qB_jK_j^\dagger$ with the
isometries $K_j$ that place the bond pairs of block $j$, the polar decomposition of
$B$ is $V=\sum_jV_jK_j^\dagger$, $P=\sum_j\mu_j^qK_jP_jK_j^\dagger$ for $\mu_j>0$.
The states of distinct blocks, before and after the replacement by the fixed
points, are then exactly orthogonal, the term $O(e^{-N/\xi_{\mathrm{off\text{-}diag}}})$
vanishes, and part~(ii) reduces to part~(i) for each block, with no condition
$q=o(N)$. For multiplicity one the phase of $\mu_j$ can be absorbed into $A_j$,
which stays normal with the same transfer map, so $\mu_j>0$ is no loss. This
repair is formalized.\footnote{Formal declarations:
    \leanid{Matrix.polarIso_sum_of_orthogonal} in
    \doclink{TNLean/MPS/Preparation/OrthogonalSumPolar};
    \leanid{MPSTensor.nonNormalApproxVector_blockSum} in
    \doclink{TNLean/MPS/Preparation/OrthogonalBlockSum};
    \leanid{MPSTensor.exists_approximationError_le_blockSum} and
    \leanid{MPSTensor.exists_approximationError_le_mul_blockSum} in
    \doclink{TNLean/MPS/Preparation/OrthogonalBlockError}.}
The orthogonality holds for instance when distinct blocks use disjoint sets of
physical indices.

The second keeps the source's hypotheses and adds to the rate a term controlled
by the mixed transfer matrices, of the form $(N/q)\,\tau^{2q}$ suggested by the
counterexample; it needs a perturbation bound for the square root of $B^\dagger B$
around its block-diagonal part. It is not formalized.

\bibliographystyle{alpha}
\bibliography{references}
\end{document}
